\documentclass[11pt,a4paper]{article}

% ============================================================
% Packages
% ============================================================

\usepackage[utf8]{inputenc}
\usepackage[T1]{fontenc}
\usepackage{lmodern}
\usepackage[english]{babel}

\usepackage{amsmath}
\usepackage{amssymb}
\usepackage{booktabs}
\usepackage{array}
\usepackage{longtable}
\usepackage{tabularx}
\usepackage{enumitem}
\usepackage{geometry}
\usepackage{xcolor}
\usepackage{listings}
\usepackage{hyperref}
\usepackage{fancyhdr}
\usepackage{titlesec}

% ============================================================
% Page Setup & Colors
% ============================================================

\geometry{
    top=2.5cm,
    bottom=2.5cm,
    left=2.5cm,
    right=2.5cm,
    headheight=14.5pt
}

\definecolor{arenaBlue}{RGB}{24, 85, 144}
\definecolor{codebg}{RGB}{247, 249, 250}
\definecolor{codegreen}{RGB}{40, 140, 60}
\definecolor{codegray}{RGB}{100, 100, 100}
\definecolor{codeblue}{RGB}{20, 70, 160}

\hypersetup{
    colorlinks=true,
    linkcolor=arenaBlue,
    urlcolor=arenaBlue,
    citecolor=arenaBlue
}

% ============================================================
% Header / Footer Configuration
% ============================================================

\pagestyle{fancy}
\fancyhf{}

\lhead{\textbf{Algorithm Arena}}
\rhead{Problem Specification: Axelrod IPD}

\cfoot{\thepage}

% ============================================================
% Custom Problem Metadata Definitions
% ============================================================

\newcommand{\problemname}{Iterated Prisoner's Dilemma}
\newcommand{\shortname}{axelrod-ipd}
\newcommand{\problemversion}{1.0.0}
\newcommand{\problemdate}{August 2026}

% ============================================================
% Document Content
% ============================================================

\begin{document}

% ------------------------------------------------------------
% Title Page
% ------------------------------------------------------------
\hypersetup{pageanchor=false}
\begin{titlepage}

    \centering
    \vspace*{1.0cm}

    {\Large\textbf{\color{arenaBlue}ALGORITHM ARENA}}\\[0.4cm]
    {\color{codegray!60}\rule{0.5\textwidth}{0.8pt}}\\[1.2cm]

    {\fontsize{26}{34}\selectfont\textbf{Iterated Prisoner's Dilemma}}\\[0.4cm]
    {\Large\color{arenaBlue}\textbf{Axelrod's Classic 2-Player Game}}\\[1.2cm]

    {\large\textsc{Document A: Formal Problem Specification}}\\[1.5cm]

    \begin{tabular}{ll}
        \toprule
        \textbf{Problem Identifier:} & \shortname \\
        \textbf{Specification Level:}& Document A (Pure Problem Definition) \\
        \textbf{Version:}            & \problemversion \\
        \textbf{Category:}           & Game Theory / Multi-Agent Decision Making \\
        \textbf{Difficulty:}         & Intermediate \\
        \textbf{Date:}               & \problemdate \\
        \bottomrule
    \end{tabular}

    \vfill

    \parbox{0.85\textwidth}{
        \centering\textit{
            This document establishes the mathematical and domain-level specification
            of the Iterated Prisoner's Dilemma benchmark, independent of platform
            execution runtimes and tournament matchmaking implementations.
        }
    }

\end{titlepage}
\hypersetup{pageanchor=true}

% ------------------------------------------------------------
% Table of Contents
% ------------------------------------------------------------
\pagenumbering{roman}
\tableofcontents
\newpage
\pagenumbering{arabic}

% ============================================================
% 1. Metadata
% ============================================================

\section{Metadata}

\subsection{Problem Name}
\begin{quote}
\textbf{Iterated Prisoner's Dilemma (Axelrod IPD)}
\end{quote}

\subsection{Short Identifier}
\begin{quote}
\texttt{axelrod-ipd}
\end{quote}

\subsection{Category}
\begin{itemize}
    \item[$\blacksquare$] Game Theory
    \item[$\blacksquare$] Multi-Agent Systems
    \item[$\blacksquare$] Online Decision Making
    \item[$\square$] Continuous Control
    \item[$\square$] Combinatorial Optimization
\end{itemize}

\subsection{Difficulty Level}
\begin{itemize}
    \item[$\square$] Beginner
    \item[$\blacksquare$] Intermediate
    \item[$\square$] Advanced
    \item[$\square$] Research
\end{itemize}

\subsection{Version History}

\begin{center}
\begin{tabularx}{\textwidth}{llXX}
\toprule
\textbf{Version} & \textbf{Date} & \textbf{Author} & \textbf{Changelog Summary} \\
\midrule
0.1.0 & 2026-08-19 & Arena Core Team & Initial concept formulation and problem outline \\
1.0.0 & 2026-08-20 & Arena Core Team & Complete standardized problem specification (Document A) \\
\bottomrule
\end{tabularx}
\end{center}

% ============================================================
% 2. Problem Overview
% ============================================================

\section{Problem Overview}

\subsection{Description}
The \textit{Prisoner's Dilemma} is the canonical non-zero-sum $2 \times 2$ symmetric game in modern mathematical game theory. Two agents interact by choosing simultaneously between two discrete actions: \textbf{Cooperation} ($C$) and \textbf{Defection} ($D$). In a single-shot interaction, defection strictly dominates cooperation for each individual agent regardless of the opponent's choice; however, mutual cooperation results in a strictly higher joint welfare than mutual defection.

In the \textbf{Iterated Prisoner's Dilemma (IPD)}, the stage game is repeated sequentially over a fixed horizon of $T = 200$ rounds. Because players maintain complete memory of past joint decisions, repetition unlocks complex reciprocal dynamics, enabling retaliation against defection, mutual trust building, punishment, and forgiveness.

\subsection{Motivation}
The Iterated Prisoner's Dilemma serves as the inaugural benchmark for Algorithm Arena because it exemplifies fundamental algorithmic principles:
\begin{enumerate}
    \item \textbf{Tension Between Rationality and Welfare:} Captures the mathematical conflict between myopic personal payoff maximization and long-term reciprocal cooperation.
    \item \textbf{Opponent-Dependent Dynamics:} The effectiveness of any decision rule cannot be evaluated in isolation; its performance is strictly coupled to the behavior and memory of the opponent.
    \item \textbf{Clean Strategic Depth:} Despite minimal action space ($|\mathcal{A}| = 2$), optimal behavior requires state estimation, opponent profiling, and adaptive retaliation.
\end{enumerate}

\subsection{Learning Objectives}
By analyzing and designing algorithms for this problem, researchers and competitors will master:
\begin{itemize}
    \item The formalization of stage games versus repeated games with finite horizons.
    \item Nash equilibrium versus Subgame-Perfect Equilibrium (SPE) and the Folk Theorem.
    \item Designing reactive, memory-bounded, and stochastic policies in adversarial and cooperative multi-agent environments.
\end{itemize}

% ============================================================
% 3. Formal Problem Definition
% ============================================================

\section{Formal Problem Definition}

\subsection{Agents and Roles}
Each match consists of $N = 2$ symmetric autonomous agents:
\[
\text{Agent } 1 \quad \text{and} \quad \text{Agent } 2
\]
Both agents possess identical action sets, payoff structures, and observation formats.

\subsection{Horizon and Timing}
\begin{itemize}
    \item \textbf{Game Horizon:} Fixed finite duration of $T = 200$ sequential rounds: $t \in \{0, 1, \dots, T-1\}$.
    \item \textbf{Decision Timing:} Discrete-time, synchronous lock-step execution. In each round $t$, both agents select their actions simultaneously without observing the opponent's current choice $a_t^{(\text{opp})}$.
\end{itemize}

\subsection{Action Space}
At each round $t$, agent $i \in \{1, 2\}$ selects an action $a_t^{(i)}$ from the discrete action set $\mathcal{A}$:
\[
\mathcal{A} = \{C, D\} \equiv \{1, 0\}
\]
where:
\begin{itemize}
    \item $C = 1$: Cooperate
    \item $D = 0$: Defect
\end{itemize}

\subsection{Payoff Matrix}
The stage game payoffs for joint action pairs $(a^{(1)}, a^{(2)}) \in \mathcal{A} \times \mathcal{A}$ are defined by the canonical Axelrod payoff matrix:

\begin{center}
\begin{tabular}{cc|cc}
\multicolumn{2}{c}{} & \multicolumn{2}{c}{\textbf{Agent 2}} \\
& & \textbf{Cooperate ($C$)} & \textbf{Defect ($D$)} \\
\hline
\textbf{Agent 1} & \textbf{Cooperate ($C$)} & $(R, R) = (3, 3)$ & $(S, T) = (0, 5)$ \\
& \textbf{Defect ($D$)} & $(T, S) = (5, 0)$ & $(P, P) = (1, 1)$ \\
\end{tabular}
\end{center}

The payoff parameters satisfy the two mandatory mathematical conditions of the Prisoner's Dilemma:
\begin{enumerate}
    \item \textbf{Strict Incentive Order:}
    \[
    T > R > P > S \quad \iff \quad 5 > 3 > 1 > 0
    \]
    where $T$ represents Temptation, $R$ represents Reward, $P$ represents Punishment, and $S$ represents Sucker's payoff.
    \item \textbf{Social Welfare Efficiency:}
    \[
    2R > T + S \quad \iff \quad 2(3) > 5 + 0 \quad (6 > 5)
    \]
    ensuring that sustained mutual cooperation yields strictly higher aggregate payoff than alternating exploitation.
\end{enumerate}

\subsection{State and Observation Space}
Let $h_t$ represent the public history of joint actions up to round $t-1$:
\[
h_t = \Big( (a_0^{(1)}, a_0^{(2)}), (a_1^{(1)}, a_1^{(2)}), \dots, (a_{t-1}^{(1)}, a_{t-1}^{(2)}) \Big) \in (\mathcal{A} \times \mathcal{A})^t
\]
At $t = 0$, the history is empty: $h_0 = \emptyset$.

At round $t$, the observation provided to agent $i$ is:
\[
o_t^{(i)} = \Big( t, \; \mathbf{a}_{0:t-1}^{(i)}, \; \mathbf{a}_{0:t-1}^{(j)} \Big)
\]
where $\mathbf{a}_{0:t-1}^{(i)}$ is the sequence of the agent's own past actions and $\mathbf{a}_{0:t-1}^{(j)}$ is the sequence of the opponent's past actions.

\subsection{Transition Function}
The environment state transitions deterministically by appending the newly selected joint actions:
\[
h_{t+1} = h_t \cup \{(a_t^{(1)}, a_t^{(2)})\}
\]

\subsection{Reward Function and Total Match Utility}
The stage reward $r_t^{(i)}$ for agent $i$ against opponent $j$ is:
\[
r_t^{(i)} = R_i(a_t^{(i)}, a_t^{(j)}) =
\begin{cases}
R = 3, & \text{if } a_t^{(i)} = 1 \text{ and } a_t^{(j)} = 1 \\
S = 0, & \text{if } a_t^{(i)} = 1 \text{ and } a_t^{(j)} = 0 \\
T = 5, & \text{if } a_t^{(i)} = 0 \text{ and } a_t^{(j)} = 1 \\
P = 1, & \text{if } a_t^{(i)} = 0 \text{ and } a_t^{(j)} = 0
\end{cases}
\]

The cumulative utility for agent $i$ in a match of length $T = 200$ is:
\[
U_i(h_T) = \sum_{t=0}^{T-1} r_t^{(i)}(a_t^{(i)}, a_t^{(j)})
\]

% ============================================================
% 4. Mathematical Model
% ============================================================

\section{Mathematical Model}

\subsection{Notation Summary}

\begin{tabularx}{\textwidth}{llX}
\toprule
\textbf{Symbol} & \textbf{Domain} & \textbf{Meaning} \\
\midrule
$N$ & $\{2\}$ & Number of agents in a match \\
$t$ & $\{0, 1, \dots, T-1\}$ & Round index \\
$T$ & $\mathbb{N}$ ($200$) & Total rounds per match \\
$\mathcal{A}$ & $\{0, 1\} \equiv \{D, C\}$ & Discrete action space \\
$a_t^{(i)}$ & $\mathcal{A}$ & Action chosen by agent $i$ at round $t$ \\
$h_t$ & $(\mathcal{A} \times \mathcal{A})^t$ & History sequence of joint actions up to round $t-1$ \\
$r_t^{(i)}$ & $\{0, 1, 3, 5\}$ & Stage payoff received by agent $i$ at round $t$ \\
$U_i$ & $[0, 5T] \equiv [0, 1000]$ & Cumulative match payoff of agent $i$ \\
$\pi_i$ & $\mathcal{H} \to \mathcal{P}(\mathcal{A})$ & Decision policy mapping observable history to action \\
\bottomrule
\end{tabularx}

\subsection{Theoretical Properties}
\begin{enumerate}
    \item \textbf{Single-Stage Dominance:} Defection strictly dominates cooperation:
    \[
    \forall a^{(2)} \in \{0, 1\}, \quad R_1(0, a^{(2)}) > R_1(1, a^{(2)})
    \]
    The unique pure-strategy Nash equilibrium of the single-shot game is $(D, D) = (0, 0)$, producing $(1, 1)$.
    \item \textbf{Folk Theorem in Repeated Settings:} In an infinite or unknown horizon game with discount factor $\delta \ge \frac{T - R}{T - P} = 0.5$, mutual cooperation is supported as a Subgame-Perfect Equilibrium through trigger retaliation. In known fixed finite games of length $T$, backward induction predicts mutual defection in theoretical self-play; however, in empirical multi-agent tournaments with boundedly rational opponents, reciprocal cooperative strategies dominate.
\end{enumerate}

% ============================================================
% 5. Concrete Match Walkthroughs
% ============================================================

\section{Concrete Match Examples}

\subsection{Example 1: Mutual Cooperation (Tit-for-Tat vs. Pavlov)}
\begin{center}
\small
\begin{tabularx}{\textwidth}{cccccX}
\toprule
\textbf{Round} $t$ & \textbf{Agent 1 (TFT)} & \textbf{Agent 2 (WSLS)} & \textbf{Payoff 1} & \textbf{Payoff 2} & \textbf{Outcome} \\
\midrule
$0$ & $C$ ($1$) & $C$ ($1$) & $3$ & $3$ & Mutual Cooperation \\
$1$ & $C$ ($1$) & $C$ ($1$) & $3$ & $3$ & Mutual Cooperation \\
$2$ & $C$ ($1$) & $C$ ($1$) & $3$ & $3$ & Mutual Cooperation \\
$3$ & $C$ ($1$) & $C$ ($1$) & $3$ & $3$ & Mutual Cooperation \\
$4$ & $C$ ($1$) & $C$ ($1$) & $3$ & $3$ & Mutual Cooperation \\
\midrule
\textbf{Total (5 rds)} & & & \textbf{15} & \textbf{15} & \textbf{Tie (Social Welfare Optimal)} \\
\bottomrule
\end{tabularx}
\end{center}

\subsection{Example 2: Exploitation & Retaliation (Tit-for-Tat vs. Always Defect)}
\begin{center}
\small
\begin{tabularx}{\textwidth}{cccccX}
\toprule
\textbf{Round} $t$ & \textbf{Agent 1 (TFT)} & \textbf{Agent 2 (ALLD)} & \textbf{Payoff 1} & \textbf{Payoff 2} & \textbf{Outcome} \\
\midrule
$0$ & $C$ ($1$) & $D$ ($0$) & $0$ & $5$ & Agent 2 Exploits Sucker \\
$1$ & $D$ ($0$) & $D$ ($0$) & $1$ & $1$ & Agent 1 Retaliates \\
$2$ & $D$ ($0$) & $D$ ($0$) & $1$ & $1$ & Mutual Defection \\
$3$ & $D$ ($0$) & $D$ ($0$) & $1$ & $1$ & Mutual Defection \\
$4$ & $D$ ($0$) & $D$ ($0$) & $1$ & $1$ & Mutual Defection \\
\midrule
\textbf{Total (5 rds)} & & & \textbf{4} & \textbf{9} & \textbf{Agent 2 Wins Match} \\
\bottomrule
\end{tabularx}
\end{center}

% ============================================================
% 6. References
% ============================================================

\section{References}

\begin{enumerate}
    \item Axelrod, R. (1980). \textit{Effective Choice in the Prisoner's Dilemma}. Journal of Conflict Resolution, 24(1), 3-25.
    \item Axelrod, R. (1984). \textit{The Evolution of Cooperation}. Basic Books, New York.
    \item Nowak, M. A., \& Sigmund, K. (1993). \textit{A strategy of win-stay, lose-shift that outperforms tit-for-tat in the Prisoner's Dilemma game}. Nature, 364(6432), 56-58.
    \item Press, W. H., \& Dyson, F. J. (2012). \textit{Iterated Prisoner's Dilemma contains strategies that dominate any evolutionary opponent}. Proceedings of the National Academy of Sciences, 109(26), 10409-10413.
\end{enumerate}

\end{document}
